\section{从 VAE 到层级化 VAE}

\subsection{增加层级的动机}

VAE 的局限在于单一高斯分布难以拟合复杂的数据分布。一个自然的思路是：不再使用单一的潜变量 $z$，而是引入\textbf{多层潜变量} $x_1, x_2, \ldots, x_T$，形成层级化结构。

\begin{knowledgebox}{层级化提升表达能力的直觉}
单个高斯分布只有均值和方差两组参数，表达能力有限。但如果我们将多个高斯转移分布的乘积作为整体的似然分布，那么虽然每一步的转移仍是简单的高斯，\textbf{整体分布可以远比单个高斯复杂}。这就像一条简单的直线路径可以通过多次小的方向调整来逼近任意复杂的曲线。
\end{knowledgebox}

\subsection{符号切换}

从这里开始，我们切换符号体系以适应扩散模型的框架：

\begin{itemize}
    \item $x_0$：原始数据点（之前的 $x$）
    \item $x_1, x_2, \ldots, x_T$：多层潜变量（之前的 $z$）
    \item $x_T$：最终的纯噪声，服从 $\mathcal{N}(0, I)$
\end{itemize}

\begin{warningbox}{符号变化提醒}
注意从 VAE 到扩散模型的符号切换：之前用 $z$ 表示潜变量、$x$ 表示数据，现在统一用 $x_t$ 表示不同时间步的变量。这是因为在扩散模型中，所有中间变量与数据具有\textbf{相同的维度}，它们处于同一空间中。不同论文可能使用不同的符号体系，阅读时需要注意。
\end{warningbox}

\subsection{马尔可夫假设}

为了简化层级 VAE 的训练，我们引入\textbf{马尔可夫假设}：每个 $x_t$ 的采样仅依赖于前一步 $x_{t-1}$，而不依赖于更早的状态。

$$q(x_t | x_{t-1}, x_{t-2}, \ldots, x_0) = q(x_t | x_{t-1})$$

在马尔可夫假设下，整个联合分布可以分解为转移分布的乘积：

\textbf{生成过程}（似然/逆过程）：
$$p_\theta(x_0, x_1, \ldots, x_T) = p(x_T) \prod_{t=1}^{T} p_\theta(x_{t-1} | x_t)$$

\textbf{编码过程}（变分后验/前向过程）：
$$q(x_1, \ldots, x_T | x_0) = \prod_{t=1}^{T} q(x_t | x_{t-1})$$

\begin{importantbox}{马尔可夫假设的重要性}
如果不假设马尔可夫性质，$x_t$ 的采样可能依赖于 $x_0, x_1, \ldots, x_{t-1}$ 的所有历史状态，这将使训练变得极其复杂。马尔可夫假设极大地简化了计算，使得 ELBO 可以被有效分解和优化。后续课程将讨论非马尔可夫的情况（如 DDIM）。
\end{importantbox}

\subsection{层级 VAE 的挑战}

虽然层级结构增加了表达能力，但在标准 VAE 框架下仍面临挑战：

\begin{enumerate}
    \item \textbf{联合训练困难}：需要同时训练编码器（学习变分后验 $q_\phi$）和解码器（学习似然 $p_\theta$），层级越多训练越不稳定。
    \item \textbf{计算代价高}：ELBO 中的一致性项（consistency term）涉及多个随机变量，计算昂贵。
\end{enumerate}

这些挑战催生了扩散模型的核心设计选择：\textbf{固定前向过程，只学习逆过程}。

\subsection{本章小结}

通过引入多层潜变量和马尔可夫假设，层级化 VAE 增强了表达能力。然而，联合训练编码器和解码器仍然困难，这促使我们寻找一种无需学习编码器的方案——即扩散模型。


